\section{Introduction}

Durian (\textit{Durio zibethinus}), often referred to as the ``King of Fruits,'' is a high-value tropical crop cultivated extensively across Southeast Asia, particularly in Thailand, Malaysia, Indonesia, and Vietnam. The global durian market has experienced remarkable growth, with export values exceeding \$3 billion annually, driven by increasing demand from China and other international markets. However, durian cultivation faces significant challenges from various foliar diseases, including Algal Leaf Spot (\textit{Cephaleuros virescens}), Leaf Blight (\textit{Phytophthora palmivora}), and Leaf Spot (\textit{Colletotrichum} spp.), which can reduce photosynthetic capacity, weaken trees, and ultimately decrease fruit yield and quality by up to 30--40\% if left unmanaged \cite{50}, \cite{51}.

Traditional disease management in durian orchards relies on manual visual inspection by trained agronomists or experienced farmers, a process that is labor-intensive, time-consuming, subjective, and prone to inconsistencies. Early-stage disease symptoms are often subtle and easily overlooked, leading to delayed intervention and widespread infection. Furthermore, the increasing scale of commercial durian plantations (often exceeding hundreds of hectares) makes comprehensive manual monitoring impractical. These limitations underscore the urgent need for automated, accurate, and scalable disease detection systems to support precision agriculture practices \cite{52}, \cite{53}.

Recent advances in deep learning, particularly convolutional neural networks (CNNs) for computer vision, have demonstrated remarkable success in automated plant disease detection across various crops. Among CNN architectures, the You Only Look Once (YOLO) family of object detectors has emerged as the dominant paradigm for real-time detection tasks due to its superior balance between accuracy and inference speed \cite{17}, \cite{21}, \cite{29}. The YOLO series has undergone continuous evolution from YOLOv1 (2015) to the latest YOLOv12 (2024), with each iteration introducing architectural innovations aimed at improving detection accuracy, reducing computational complexity, and enhancing deployment flexibility \cite{30}, \cite{31}, \cite{32}.

While GPU-accelerated cloud-based inference offers high accuracy and throughput, it introduces latency, bandwidth costs, privacy concerns, and dependency on stable internet connectivity—factors that are particularly problematic in rural agricultural settings. Edge computing, which performs inference directly on low-power devices deployed in the field, addresses these limitations by enabling real-time, offline, and privacy-preserving disease detection \cite{44}, \cite{47}, \cite{49}. The Raspberry Pi, a low-cost (approximately \$35--\$75), credit-card-sized single-board computer, has gained widespread adoption in agricultural IoT applications due to its affordability, versatility, and active community support. However, its limited computational resources (ARM Cortex-A72 CPU, 4--8 GB RAM, no dedicated GPU) pose significant challenges for deploying state-of-the-art deep learning models \cite{22}, \cite{42}.

Despite extensive research on YOLO-based plant disease detection and edge deployment strategies, several critical gaps remain. First, most existing studies focus on older YOLO versions (YOLOv3--YOLOv5) or evaluate only a single model variant, lacking systematic comparisons of the latest architectures (YOLOv8--YOLOv12) that incorporate recent innovations such as C2f modules, decoupled heads, and anchor-free detection \cite{4}, \cite{14}, \cite{25}. Second, while numerous studies report GPU-based performance metrics, comprehensive evaluations on resource-constrained edge devices like Raspberry Pi remain scarce, particularly for durian disease detection \cite{19}, \cite{48}. Third, the impact of different model optimization formats (PyTorch, ONNX, TensorFlow Lite, OpenVINO) on edge inference performance is often overlooked, despite their critical importance for practical deployment \cite{47}.

This study addresses these gaps by presenting the first systematic comparison of next-generation YOLO variants (YOLOv8, YOLOv10, YOLOv11, YOLOv12) for real-time durian leaf disease detection on Raspberry Pi edge devices. Our comprehensive evaluation encompasses both nano (n) and small (s) model series across multiple optimization formats, providing actionable insights for practitioners seeking to deploy YOLO-based disease detection systems in resource-constrained agricultural environments.

\subsection{Research Objectives and Contributions}

The primary objectives of this research are:

\begin{enumerate}
\item \textbf{Systematic Benchmarking}: Conduct a comprehensive performance comparison of YOLOv8, YOLOv10, YOLOv11, and YOLOv12 (nano and small variants) for durian leaf disease detection, evaluating both accuracy metrics (mAP@50, mAP@50:95, Precision, Recall, F1-score) and computational efficiency (inference time, FPS, model size, memory footprint).

\item \textbf{Edge Deployment Evaluation}: Assess the feasibility and performance of deploying these models on Raspberry Pi 4B (4 GB RAM), a representative low-cost edge device, across multiple optimization formats (PyTorch, ONNX, TFLite, OpenVINO).

\item \textbf{Optimization Format Analysis}: Quantify the impact of different model optimization and quantization strategies on inference speed, accuracy retention, and resource utilization to identify optimal deployment configurations.

\item \textbf{Practical Guidelines}: Provide evidence-based recommendations for selecting appropriate YOLO variants and optimization strategies for real-world agricultural edge computing applications, considering accuracy-speed-resource trade-offs.

\item \textbf{Benchmark Dataset}: Introduce a curated durian leaf disease dataset with approximately 2,000 annotated images covering four classes (Algal Leaf Spot, Leaf Blight, Leaf Spot, Healthy), which will be made publicly available to support future research.
\end{enumerate}

The key contributions of this work include:

\begin{itemize}
\item \textbf{First Comprehensive Comparison}: To our knowledge, this is the first study to systematically compare YOLOv8, YOLOv10, YOLOv11, and YOLOv12 for plant disease detection on edge devices, providing critical insights into the practical implications of recent YOLO architectural innovations.

\item \textbf{Multi-Format Optimization Analysis}: We provide the first detailed evaluation of PyTorch, ONNX, TFLite, and OpenVINO optimization formats for YOLO models on Raspberry Pi, revealing significant performance differences (up to 3$\times$ speedup) that are crucial for deployment decisions.

\item \textbf{Durian Disease Detection Benchmark}: We establish baseline performance metrics for durian leaf disease detection using state-of-the-art YOLO models, addressing a significant gap in tropical fruit disease detection research.

\item \textbf{Actionable Deployment Guidelines}: Our findings demonstrate that YOLOv11n with ONNX optimization achieves the optimal balance (87.3\% mAP@50, 12.5 FPS) for real-time edge deployment, providing a practical reference point for agricultural AI practitioners.

\item \textbf{Open-Source Resources}: We will release our dataset, trained model weights, and inference code to facilitate reproducibility and accelerate research in agricultural edge AI.
\end{itemize}

The remainder of this paper is organized as follows: Section II reviews related work on YOLO evolution, lightweight architectures, edge deployment, and durian disease detection. Section III describes our dataset, experimental methodology, and evaluation metrics. Section IV presents comprehensive results comparing model accuracy, inference performance, and resource utilization. Section V discusses key findings, practical implications, limitations, and future research directions. Section VI concludes the paper with a summary of contributions and recommendations.
